%% kapa: exam=1 round=32 year=2021 period=1 form=A subject=real_estate subject_ko=부동산학원론 seq=10 number=90
%% source: Q-Net 감정평가사 기출문제 (cst00309) artlSeq=5208397 file=extracted/1008/5208397/A형/2021년도 제32회 감정평가사 1차 1교시 A형.pdf
%% conversion: kordoc 4.24.51 -> tools/gen1.py (auto); review=proofread-sonnet; flags=none
\begin{question}{90}
다음 중 현행 부동산가격공시제도에 관한 설명으로 옳은 것은 몇 개인가?
\begin{framedtext}
○ 표준주택가격의 조사·평가는 감정평가사가 담당한다. \par
○ 개별주택가격의 공시기준일이 6월 1일인 경우도 있다. \par
○ 공동주택가격의 공시권자는 시장·군수·구청장이다. \par
○ 표준지공시지가는 표준지의 사용·수익을 제한하는 사법상의 권리가 설정되어 있는 경우 이를 반영하여 평가한다. \par
○ 개별공시지가는 감정평가법인 등이 개별적으로 토지를 감정평가하는 경우에 기준이 된다. \par
\end{framedtext}
\begin{choices}
  \item 없음
  \item 1개
  \item 2개
  \item 3개
  \item 4개
\end{choices}
\end{question}
\answer{2}
